\chapter{Introduction}

Les automates cellulaires~(AC)\footnote{%
  Pour des collections thématiques voir
  \cite{Demongeot85,Wolfram86a,Legendi87,Gutowitz90a}, %
  pour des ouvrages thématiques voir
  \cite{Preston84,Toffoli87,Jackson91c}, %
  pour une approche didactique voir %
  \cite{Casti89a,Schroeder90a}.} %
sont des systèmes dynamiques virtuels, %
discrets en espace, en temps et en états, homogènes, synchrones et
déterministes. %
L'espace est constitué de cellules identiques dotées d'un réseau
d'interconnexion régulier~(une grille de dimension~2 ou~3 par
exemple). %
Une cellule est une variable discrète~(une chaîne de bits dont la
taille détermine le nombre d'états de la cellule). %
Le temps est une séquence d'espaces générée par l'application
synchrone d'une fonction de transition déterministe. %
La fonction de transition est locale dans le temps et dans l'espace,
son argument est un ensemble fini d'états de cellules voisines du
réseau d'interconnexion, pour un nombre fini de points antérieurs. %

Les AC sont sujets d'études ou d'applications dans les domaines des
mathématiques %
\cite[]{Jen90,Vichniac90}, %
de l'architecture des ordinateurs %
\cite[]{Legendi87a,Zsoter87}, %
de la programmation parallèle %
\cite[]{Preston84,Almasi87} %
et de la simulation de systèmes naturels %
tant physiques\footnote{%
  En particulier les \aquote{Lattice Gas Automata}~\cite[]{Despain90,Margolus90}.} %
\cite[]{Margolus84,Toffoli84,Rossi93} %
que biologiques\footnote{%
  L'étude de la vie artificielle inspire de nombreux travaux sur les
  AC, souvent assez éloignés d'observations biologiques. %
  Une exception notable étant \cite{Hameroff89}, qui propose
  l'existence d'une capacité au calcul cellulaire pour le
  cytosquelette.} %
\cite[]{Tamayo89,Boer93}. %
Ils constituent un champ d'études des mathématiques discrètes et un
modèle important en théorie du calcul\footnote{%
  Voir par exemple \cite{Lindgren90} pour une équivalence
  entre certains AC~1D et une machine de Turing.}, %
en particulier pour les problèmes liés au parallélisme\footnote{%
  Les AC mathématiques mettent en jeu
  un espace infini, en particulier à limites périodiques. %
  Notre point de vue est également empirique et
  implémentatoire: nous viserons spécifiquement la réalisation de
  programmes dotés du meilleur rendement. %
  En conséquence, dans un contexte de ressources mémoire limitées, nos
  AC seront évidemment dotés d'un espace restreint, par défaut torique.}; %
leur structure est très proche de l'architecture des machines
SIMD\footnote{%
  \cite{Greussay85} indique que la structure des AC est très proche de
  l'architecture des machines SIMD, et nous pensons que le progrès des
  connaissances et des outils dans ces deux domaines sont intimement
  liés. %
  Les travaux de \cite{Rozin87,Rozin88,Rozin93} sur la programmation
  plane et de \cite{Signorini92,signorini94} sur la programmation par
  configurations cellulaires témoignent de la fertilité de cette
  approche. %
  La puissance de calcul et la taille de la mémoire disponible sur des
  machines aujourd'hui largement accessibles permet d'envisager la
  réalisation de multiples expériences dans le domaine du parallélisme
  synchrone massif, sans attendre l'arrivée de machines SIMD\@. %
  Les AC fournissent un outil conceptuel en mesure de jouer un rôle
  central dans la recherche d'objets naturels~(comportement émergent)
  ou artificiels~(comportement construit), pouvant servir de base à la
  définition d'outils pratiques d'exploitation des ressources du
  parallélisme synchrone massif.}; %
indépendamment de l'architecture des machines, de nombreux
algorithmes se prêtent à une formulation
cellulaire~\cite[]{Uhr87,Uhr84}, en particulier les algorithmes de
traitement d'images~\cite[]{Preston84,Hasselbring92}; enfin, les AC
peuvent être vus comme une version discrète de systèmes
d'équations différentielles\footnote{%
  Et à ce titre permettent la modélisation de nombreux systèmes
  naturels. %
  La modélisation d'un système naturel par un ensemble d'équations
  différentielles établit des relations différentielles entre les
  variables macroscopiques du système~(des populations par exemple). %
  La modélisation par AC établit des relations entre les composants
  discrets élémentaires~(les cellules) du système. %
  Les variables macroscopiques deviennent des propriétés émergentes du
  système dont le choix peut être fait a posteriori, mais le modèle
  doit impérativement reposer sur une description en termes de
  comportement collectif du phénomène étudié. %

  La description d'un système naturel en termes de comportement
  collectif n'implique pas l'identification d'une cellule à un
  individu. %
  Par exemple, dans une simulation de dynamique des populations, une
  cellule peut être identifiée à une densité de population en un point
  d'espace. %
  Néanmoins la validité de la modélisation~(conformité de la dynamique
  des populations d'un AC à un système d'équations différentielles
  standard de référence par exemple) dépend toujours de la présence
  d'un nombre minimum de cellules dans la simulation.}. %

\section[Objectifs et contraintes]{Nos objectifs et contraintes pour un
  environnement de dé\-ve\-lop\-pe\-ment et d'observation d'AC}

Nous cherchons à définir et à réaliser un environnement de
développement et d'observation de ces systèmes. %
Une des motivations principales de ce travail réside dans le désir de
voir~(et montrer) les séquences animées générées par un AC\@. %
En effet s'il est aisé de construire des AC en vue de donner des
instantanés photographiques d'états successifs, il est plus difficile de
les visualiser en continu filmique. %
Bien qu'il existe des outils\footnote{%
  \mysoftcite[Cellsim]{cellsim}, %
  \mysoftcite{scamper}, %
  \mysoftcite[CA~LAB]{ca_lab}~\cite[]{Rucker89}, %
  \mysoftcite{xlife}, %
  \mysoftcite[HODGE-C]{hodge_c}, %
  \mysoftcite{cellang}. %
  Les numéros en indice supérieur entre crochets renvoient au répertoire des
  outils, donné en appendice.} %
de spécification d'AC et d'observation de leur dynamique, aucun de
ceux dont nous connaissons l'existence ne répond
pleinement aux contraintes que nous allons définir. %
Nous souhaitons également que cet outil permette immédiatement la
recherche et la création de nouveaux AC, puis la constitution d'un
catalogue d'AC connus. %
Enfin nous espérons qu'il serve à l'identification d'AC pouvant servir
de primitives dans un futur langage basé sur la combinaison d'AC. %
Les qualités recherchées sont la rapidité, l'extensibilité, ainsi que la
facilité de spécification et d'interaction\footnote{%
  La \aconcept{rapidité} peut être définie comme le rendement d'une
  implémentation   pour une architecture donnée. %
  Dans ce contexte, et pour une architecture cible séquentielle, la
  rapidité s'exprime comme le nombre d'instructions
  élémentaires de la machine cible utilisée pour simuler une
  instruction de la machine source. %
  Cette définition est néanmoins trop générale, en effet, si nous
  proposons dans le cadre de ce travail des moteurs cellulaires à
  rendement constant~(\ie indépendant de l'espace initial et de la
  règle de transition), il est possible de concevoir des moteurs
  cellulaires à rendement variable: \mysoftcite{xlife} découpe
  l'espace cellulaire en blocs et permet le calcul sur des
  configurations cellulaires virtuelles de grandes tailles;
  \cite{Gosper84} décrit un algorithme basé sur une mémoire cache
  asynchrone des configurations cellulaires~(le cache contient un
  treillis de morceaux d'espace-temps cellulaire). %

  L'extensibilité peut, dans notre contexte, se ramener à la
  conception d'un ensemble d'objets constituant à la fois la boîte à
  outils et les éléments de base de la construction des expériences
  cellulaires. %
  La recherche de cette qualité peut entrer en conflit avec la
  recherche de la rapidité, par exemple les variations de paramètres
  sur les moteurs cellulaires à rendement constant que nous avons
  réalisés passent par une phase de macro-génération de code source.}. %


\subsection*{Le cadre système, matériel, et logiciel:~les performances}

Notre environnement de développement est Unix et X11, notre
architecture cible est une machine standard de~10~à~100~mips
et~8~à~64 \astrangeterm{mégabytes}{%
  Nous emploierons indifféremment les termes octet et byte, ainsi que
  leurs dérivés: kilobyte, mégaoctet ou gigaoctet.} %
de mémoire utile. %
Pour un AC bidimensionnel composé de \asquare{256}
cellules de~2~bits\footnote{%
  Avec un voisinage de Moore~(9~bits de voisinage), 2~bits d'états
  donnent une table de taille 512~kilobit~($2^{2 \times  9} \times  2$), ce qui
  est très loin des limites que nous nous sommes fixés, mais le passage
  à trois bits de voisinage donne une table de taille
  48~mégaoctets~($\frac{2^{3 \times  9} \times  3}{8}$) qui se trouve à la limite
  des ressources mémoires considérées. %
  Nous insistons cependant sur le fait que rien n'interdit en pratique
  l'usage de telles tables, même si la connaissance de la
  distribution des configurations locales devient nécessaire pour
  prédire les performances du système constitué du triplet
  \emph{table} de transition, \emph{configuration} initiale,
  \emph{performance} intrinsèque du système de gestion mémoire de la
  machine supportant l'application.} %
avec un voisinage de Moore, le système doit être capable de calculer
une dizaine d'états par seconde~(10~Hz)\footnote{%
  D'un point de vue \aquote{esthétique}, cette fréquence correspond à
  un minimum subjectif permettant de restituer l'illusion du mouvement
  continu~(les images télévisuelles ont une fréquence de~25~Hz, les
  images cinématographiques modernes une fréquence de~24~Hz, les
  images cinématographiques originelles une fréquence de~14 à~20~Hz,
  le passage à des fréquences plus élevées est dicté par des
  considérations techniques liées aux besoins du parlant pour le cinéma
  et à la fréquence du secteur alternatif pour la télévision). %
  Cette fréquence est donc élevée, car d'un point de vue
  \aquote{système} elle correspond~(pour notre espace de référence
  $256^2$) à un débit de~640~kilo \bialt{opérations}{cellules} par
  seconde, ce qui pour une machine hypothétique à~20~mips, correspond
  à~32~instructions de la machine cible pour une instruction de la
  machine source~(une mise à jour cellulaire).}. %
Le coût lié à la visualisation est plus difficile à évaluer, mais la
chaîne de traitement complète~(production, extraction, visualisation)
doit permettre l'observation d'un film à une fréquence supérieure
à~1~Hz\footnote{%
  Eisenstein, Gance ou Godard par leur utilisation occasionnelle de
  l'alternance de plans à une telle fréquence illustrent le procédé
  esthétique. %

  Cette fréquence correspond par ailleurs à un impératif ergonomique
  subjectif concernant le \aquote{temps de réponse} d'une application,
  elle borne la taille des espaces accessibles à un mode de
  visualisation séquentielle lorsque la contrainte de l'illusion du
  mouvement est levée. %
  
  La perte de l'illusion du mouvement, n'est pas contradictoire à
  l'étude du mouvement. %
  Le physiologiste Marey étudie le vol des oiseaux:~son %
  \aquote{fusil photographique}~(1882) et son %
  \aquote{Chronophotographe à pellicule}~(1888) %
  servent à l'analyse photographique du mouvement. %
  Bien que considéré comme précurseur du cinématographe,
  il ne se préoccupe pas de restituer l'illusion du mouvement. %
  En cinéma d'animation, la \anewterm{pixillation} utilise sciemment
  la séquence contre l'illusion. %
  \begin{quote}
    ``L'animation n'est pas l'art des dessins qui bougent, mais l'art des
    mouvements dessinés. %
    Ce qui ce passe entre chaque image a beaucoup plus d'importance que
    ce que l'on voit sur chaque image.'' \textsc{N. McLaren}.{\hfill$\cdots/\cdots$}%
  \end{quote}
  En 1941 dans ``L'évolution créatrice'', Bergson
  défend l'idée que ``le mécanisme de notre connaissance usuelle est
  de nature cinématographique''~\cite[page 305]{Bergson89}. %
  Pour lui ``la forme n'est qu'un instantané pris sur une
  transition'', seul est réel le changement continuel de forme. %
  Il place cette métaphore au cœur d'une réflexion historique qui
  remonte à Zénon d'Élée. %

  Dans son acharnement à défendre ``L'antériorité ontologique du
  continu sur le discret'', \cite{Thom92} va jusqu'à proposer
  une ``version informatisée de la preuve cartésienne de
  l'existence de Dieu'' déclarant contradictoire la finitude du nombre
  de neurones cérébraux et la pensée de l'infini. %
  Sans prétendre trancher une question aussi débattue, nous pouvons
  néanmoins suggérer que l'étude des AC propose une approche à rebours
  de cette question, la recherche pratique de l'immanence
  phénoménologique du continu remplaçant les tentatives de
  démonstration de son antériorité ontologique. %

  Sur notre configuration de développement~(un~486DX2 avec~32~mégaoctets
  de mémoire, le système d'exploitation Linux et X11 pour la
  visualisation) la fréquence typique pour une chaîne comprenant la
  génération sur un voisinage de Moore~(toujours pour un espace de
  référence $256^2$) et l'extraction d'un plan mémoire~(constitué d'un
  des deux bits d'état) dans un premier processus, pipé sur un processus
  dumpant les plans cellulaires vers le serveur X11~(tournant sur la
  même machine) est supérieure à~10~Hz. %
  Le découplage des différentes opérations de la chaîne permet
  d'obtenir une fréquence de génération supérieure à~15~Hz~(module de
  génération redirigé sur \texttt{/dev/null}) et une fréquence de
  visualisation monochrome~(pour un seul plan de bits cellulaire)
  supérieure à~30~Hz, (mesuré via un fichier historique
  de~8~mégaoctets pour~1024~générations). %
  Ces mesures s'entendent en temps écoulé réel, pour un système
  autrement au repos, les mesures de performance de visualisation ne
  constituent évidemment pas une référence à notre travail, mais elles
  permettent de considérer, au regard des mesures de performance du
  moteur cellulaire, que celui-ci atteint le débit plafond d'un moteur
  à rendement constant, \ie la meilleure performance possible est
  toujours garantie dans le pire des
  cas~(espace cellulaire et portion d'espace-temps saturé).}. %
Nous avons ainsi réalisé un noyau d'objets permettant la construction
d'expériences, sous la forme d'une librairie C, interfacée à
\asoft{Tcl}\footnote{%
  Tool Command Language.}%
~\cite[]{Ousterhout90,Ousterhout91}. %

\subsection*{Spécification:~la description des expériences cellulaires}

La spécification concerne les réseaux d'interconnexions, les fonc\-tions
de transitions, et le séquencement des opérations applicables sur
l'espace. %
Les fonc\-tions de transitions sont écrites en~\asoft{C}, et sont dotées
d'une convention d'interface prévue pour permettre une indépendance
des paramètres à leur localisation dans la mémoire d'états d'une
cellule et la réutilisation dans des fonctions de plus haut niveau. %
Les réseaux d'interconnexions sont définis et créés en~\asoft{Tcl} à
partir des primitives de la librairie. %
Les espaces sont créés en~\asoft{Tcl}. %
Le séquencement est défini par un programme~\asoft{Tcl}. %

\subsection*{Interaction:~l'interface humaine}

L'interaction se fait soit par l'intermédiaire d'un shell~\asoft{Tcl}, soit
par une interface X11 basé sur~\asoft{Tk}. %

\subsection*{Visualisation:~les outils graphiques}

La visualisation se fait par l'intermédiaire de fenêtres X11
présentant des séquences d'espaces monochromes~(plans de bits de
l'espace cellulaire) ou couleur~(tranches de bits de l'espace
cellulaire). %
Les variables macroscopiques~(éléments statistiques) sont visualisées
par l'intermédiaire d'objets X11 spécialisés~(plotters). %

\mydumpfull[htb]{casNew}%
{Une interface humaine expérimentale pour le Workbench}%
{Une interface humaine expérimentale et son dumper}%
{Une interface X11 réalisée avec
  \protect\mysoftcite{Tk}~\cite[]{Ousterhout90,Ousterhout91}. %
  Le dumper présente une image typique de la transformation d'un espace
  initial aléatoire de densité\,\Sundemi par la
  règle~\arule{Anneal}~\cite[]{Vichniac86,Toffoli87}\footnote{%
    \mygetrule[phrase]{anneal}} %
  Le dumper présente également différents compteurs liés aux temps
  propre de la règle, du séquenceur, et au pool d'espaces de l'animation. %
}

\section{Les outils existants:~pas de standard}

Une recherche aussi exhaustive que possible des outils publics
existants nous a convaincu de la nécessité d'une utilisation critique
de certains de ces outils, en particulier dans le domaine de la
visualisation, associé à une réalisation expérimentale ad hoc
d'éléments d'architecture de calcul cellulaire. %

Les outils de développement et de visualisation d'automates
cellulaires sont peu nombreux\footnote{%
  Bien que les réalisations de \arule{Life} soit
  innombrables, et les packages cellulaires plus ou moins généraux très
  nombreux, leurs réalisations minimales posant peu de problèmes, nous ne
  prenons en compte que les packages dont il existe des annonces
  explicites.}. %
Le plus connu d'entre eux est sans doute \mysoftcite{cellsim};
\mysoftcite{scamper} et \mysoftcite[CA~LAB]{ca_lab} sont également des
outils disponibles depuis 1989. %
Il existe aussi des outils spécialisés pour un modèle d'automate
cellulaire. %
\mysoftcite{xlife} fournit une réalisation d'exécution très rapide du
jeu de la vie~\cite[]{Berlekamp82}, \mysoftcite[HODGE-C]{hodge_c} une
réalisation de la machine hodge-podge de
Gerhard~\&~Schuster~\cite[]{Dewdney88}, une classe d'AC inspirée par
les réactions chimiques autocatalytiques, telle que la réaction de
Beluso\mbox{v-Z}habotinski. %
Enfin, \mysoftcite{cellang}, d'inspiration plus architecturale,
propose un langage de description d'AC~\cite[]{Eckart92a}. %
\nocite{Eckart92b}

\section[Le calcul et la visualisation]%
{Un nouvel environnement:~exploration par le calcul et la visualisation}

Nous avons donc réalisé un nouvel environnement\footnote{%
  Ci-après désigné par le terme \anewterm{workbench}.} %
de développement et de visualisation d'automates cellulaires. %

Nous destinons ce workbench tant à l'exploration de la dynamique
des automates cellulaires~(AC), qu'à la construction et à l'évaluation
d'architectures et d'algorithmes SIMD\@. %
De plus nous proposons de considérer la réalisation d'algorithmes SIMD
sur machines SISD non pas comme un palliatif à l'absence de machines
SIMD, mais comme un mode d'utilisation particulier des ressources
disponibles sur les machines SISD modernes~(50~Mips, 32~mégaoctets). %

La programmation des calculateurs massivement parallèles réels et
virtuels pose le double problème du pouvoir d'expression et de
l'efficacité d'un langage au regard d'une architecture. %
Ce problème se pose également pour les architectures SISD, mais l'on
peut considérer qu'il concerne désormais l'adéquation des langages aux
applications plutôt qu'aux architectures. %
Optant pour une approche résolument bottom-up de la recherche des
primitives, nous proposons d'utiliser les AC comme composants
élémentaires d'une combinatoire visant à la construction d'objets
\bialt{cellulaires}{SIMD} arbitrairement complexes. %
Ces objets se divisent en deux classes épistémologiques. %
Nous les nommerons objets \aconcept{analytiques} et objets
\aconcept{synthétiques}. %
Les objets analytiques sont typiquement issus de l'application
itérative d'une fonction simple sur un espace cellulaire initial
aléatoire. %
Ils recoupent intuitivement le champ d'étude traditionnel des AC\@. %
Les objets synthétiques sont issus de la sélection d'une fonction
localement universelle, et de son application contrôlée~(globalement
paramétrée) sur un espace initial soigneusement configuré. %
Ils correspondent aux architectures des machines SIMD\@. %
De l'étude des objets analytiques nous espérons tirer des
enseignements permettant à terme de considérer les processus
morphogénétiques comme des éléments de programmation. %
De la construction d'un outillage SIMD nous attendons les moyens de
contrôle nécessaires à l'étude des objets analytiques. %
Cet objectif, certes ambitieux, repose précisément sur la réalisation
par \aconcept{lookup-tables} des fonctions de bas niveau manipulables
et combinables au sein de l'environnement. %

Notre étude est organisée autour de la présentation des structures de
données choisies pour le workbench; ces structures de données
représentent \aconcept{l'espace}, \aconcept{le temps}, et
\aconcept{les lois}:~objets élémentaires du workbench. %
Nous proposons des exemples d'opérations applicables sur ces objets. %
Nous proposons également de considérer la réalisation de fonctions par
lookup-tables comme une structure de données fondamentale, tant pour
des raisons d'efficacité, que de généralité. %

Nous considérons la visualisation comme un service primordial, qui
propose l'usage du workbench en tant que laboratoire de manipulation
d'espace cellulaire~(l'espace, mais aussi les fonctions et les
dynamiques). %
Nous présentons plusieurs exemples de visualisations principalement
associées à des objets analytiques. %
Nous montrons comment peut s'effectuer le passage d'une méthode
naturaliste à une méthode artefactuelle en proposant une réalisation
par lookup-table d'une architecture SIMD spécifique, et une méthode de
transformation d'une lookup-table quelconque en séquence de contrôle
pour cette machine. %

Constatant l'absence d'outils adaptés aux contraintes de la
visualisation des espace-temps cellulaires\footnote{%
  La visualisation de volume de données dans un format inadapté
  conduit à la saturation inutile des capacités mémoire de la machine
  d'accueil d'un calcul cellulaire.}, %
constatant également l'existence d'espace-temps cellulaires binaires a
priori toriques, nous avons conçu, réalisé et testé une bibliothèque
de manipulation d'espace\footnote{%
  Coupe, séquences de coupes sur des axes de projection variables, %
  réduction intégrale sur un axe arbitraire pondéré par la distance à
  la surface de projection pour obtenir des effets de transparence.} %
cellulaires aux fins de visualisation. %
Sur cette base, nous avons par exemple pu tester l'idée d'une
visualisation relativiste des espace-temps:~%
nous avons projeté le cône de la vitesse de la lumière\footnote{%
  La vitesse de la lumière d'un espace-temps cellulaire est le rayon
  du voisinage de la famille de règle à laquelle appartient la règle
  génératrice de l'espace-temps.} %
d'un espace-temps~2D sur un espace~2D, résumant ainsi sur un seul
espace~2D la dynamique de la règle visualisée.

\mydumpfull[p]{superglider}%
{Un glider superluminique}%
{Un glider subluminique et un glider superluminique}%
{Nous présentons dans cette figure la découverte
  d'un \aquote{glider superluminique}. %
  Cette série d'espaces de configurations cellulaires est une séquence
  de tranches d'espace-temps non alignées sur l'axe du temps. %
  Malgré la distorsion introduite par ce décalage, la règle
  génératrice de cette portion d'espace-temps est
  parfaitement identifiable:~\arule{Life}\footnote{%
    Les flèches en haut à droite des deux premières photos pointent
    une configuration cyclique de~\arule{life} reconnaissables.}. %
  A première vue, la conséquence de l'utilisation d'un axe non
  standard pour la projection de la dynamique de~\arule{Life} est
  d'ordre purement spatial\footnote{%
    Cette appréciation subjective n'est applicable qu'à l'expérience
    de la vision du \aquote{film}.}, %
  exception faite du décalage introduit par le désaxage temporel. %
  Le carré sur les photos pointe vers un objet \arule{Life} de type
  glider, doté d'une vitesse de déplacement subluminique. %
  L'observation de l'espace-temps contenant la portion présentée au
  moyen de films à haute fréquence\footnote{%
    Supérieure à la fréquence télévisuelle.} %
  nous a permis d'identifier à la vue un objet paradoxal, pointé par
  un cercle sur les photos:~un glider superluminique.}

\noindent La figure~\mygetdump{superglider} présente un des usages
possible de notre librairie de manipulation d'espaces dans le
contexte du workbench:~%
nous y présentons la découverte d'un \aconcept{glider
  superluminique}\footnote{%
  nous emploierons le terme glider~(planeur) de manière générique, \ie
  pour désigner des objets cellulaires périodiques, à un décalage spatial près.}, %
en déplacement à une vitesse apparente double de la
vitesse de la lumière. %
Nous attendions évidemment de cette expérience des aberrations
temporelles, des illusions ou mirages inex\-plicables dans le contexte
d'une projection parallèle à l'axe du temps. %
Néanmoins les premiers résultats faisaient plutôt apparaître une
relative invariance de la perception visuelle subjective associée à la
règle. %
La procédure utilisée produit un déplacement apparent des
configurations, mais les objets se comportent par ailleurs, c'est à
dire relativement à l'espace apparent pour l'expérimentateur, de
manière subluminique ou luminique usuelle. %
Ce glider, le seul que nous ayons trouvé~(dans l'espace-temps
englobant la portion présentée) était donc prévisible, puisqu'il
correspond à un glider en mouvement corrélé au front de coupe, mais
néanmoins totalement surprenant. %
Beaucoup moins prévisibles nous apparaissent les conséquences
générales d'observations exploitant systématiquement des projections
d'espace-temps non parallèles à l'axe du temps.

\myfigfullR[htb]{.75}{timepyrcub}%
{Tranches d'espace-temps et pyramides relativistes}%
{Schémas de réalisation de tranches d'espace-temps et de pyramides relativistes}%
{Le cube, sur ces schémas, représente une portion d'espace-temps pour
  un AC~2D. %
  L'axe du temps est indiqué par la flèche verticale, %
  les pointillés larges sont des arêtes de plans de coupe. %
  Le schéma de gauche illustre la procédure utilisée pour réaliser les
  séquences de coupes présentées en figure~%
  \mygetdump{superglider}, %
  \mygetdump{tsum-128x3-diag}, %
  et \mygetdump{anneal-256x3-diag}. %
  Les coupes sont calculées en \aquote{distance de Manhattan}, %
  Nous avons conçu un calcul d'intersection basé sur une
  version~3D de l'algorithme de tracé de ligne de Bresenham\footnote{%
    Nous avons modifié pour cet usage la version de Bob
    Pendleton~(1992) dérivée de~\cite{Heckbert90}.}. %
  Le schéma de droite illustre la procédure utilisée pour réaliser les
  projections relativistes\footnote{%
    La projection est construite par \aconcept{cercles de voisinage}
    remontant le temps, dont le rayon détermine l'âge, éventuellement
    associés à une translation spatiale figurée par la flèche horizontale
    pleine.} %
  présentées en figures~%
  \mygetdump{life-256x3-rel}, %
  \mygetdump{life-256x3-relmov}, %
  \mygetdump{tsum-128x3-relmovedge}, %
  }

\mydumpfull[p]{tsum-128x3-diag}%
{Une séquence de tranches d'espace-temps}%
{Une séquence de tranches d'espace-temps de la règle~\arule{Qmean}}%
{La règle~\arule{Qmean} est présentée, ainsi que les autres règles
  utilisées pour la synthèse des images réalisées
  en~\mygetref[name]{the-rules}. %
  La figure~\mygetfig{timepyrcub} illustre la procédure de réalisation
  de cette séquence. %
  D'autres visualisations associées à la règle~\arule{Qmean} sont
  présentées par les figures~%
  \mygetdump{tsum-grey}, %
  \mygetdump{tsum}, %
  \mygetdump{tsum-ray-one} %
  et \mygetdump{tsum-ray-all}.}

\mydumpfull[p]{life-256x3-rel}%
{Projection relativiste fixe}%
{Une projection relativiste pour une cellule fixe}%
{Une projection relativiste~\mygetfig[paren]{timepyrcub} d'un
  espace-temps de \arule{life} pour une cellule de référence fixe,
  réalisée sur un espace-temps~\acube{256}, pour une densité
  initiale\,\Sundemi. %
  La cellule de référence est au centre de la projection, %
  celle-ci apparaît donc au sein d'un espace \aquote{âgé},
  composé en plus grand nombre d'objets point-fixe et %
  de densité moyenne  plus faible %
  que l'espace situé à proximité des bords de la projection, %
  de densité moyenne plus élevée et peuplé d'objets plus jeunes\footnote{%
    Que l'on peut cependant qualifier d'archaïques du point de vue de
    l'origine de la projection.}. %
  Il est possible d'accentuer l'âge des objets lointains en diminuant la
  vitesse de la lumière utilisée lors de la projection. %
  Un tel découplage du rayon de voisinage accroît le problème
  d'interprétation des projections produites.}

\mydumptwothirds[htb]{life-256x3-relmov}%
{Projection relativiste mobile}%
{Une projection relativiste à cellule mo\-bile}%
{Cette projection diffère de la figure~\mygetdump{life-256x3-rel} par
  l'ajout d'une translation appliquée à la cellule de référence.}

\mydumpfull[t]{anneal-256x3-rel}%
{}%
{Projection relativiste sur les 3~axes d'un cube d'espace-temps}%
{Cette figure regroupe trois projections relativistes de la
  règle~\arule{Anneal} réalisées sur une même trajectoire, %
  pour l'axe temporel et les deux axes nord-sud et est-ouest. %
  L'origine du temps est visible dans le haut de la projection pour
  ces deux axes.}

\myfigfullR[htb]{.75}{transcub}%
{Transparence globale et par plan de bits}%
{Schéma de réalisation de réductions par transparence}%
{Comme pour la figure~\mygetfig{timepyrcub}, les cubes représentent une
  portion d'espace-temps pour un AC 2D. %
  La flèche pointillée illustre une traversée orthogonale non parallèle
  à l'axe du temps, lors d'une telle traversée on peut soit réaliser un
  film\footnote{%
    A défaut de films, incompatible avec le support papier, nous
    présentons des séquences associées à une traversée non
    orthogonale~\mygetfig[paren]{timepyrcub}.}, %
  soit une réduction intégrale~(une somme de la valeur de toutes
  les cellules rencontrées). %
  Pour les réductions intégrales destinées à la production d'images,
  la valeur de la cellule est pondérée par la distance parcourue. %
  La figure~\mygetdump{tsum-128x3-trans-byte} présente une telle
  image. %
  Pour les transparences non orthogonales, les effets de bord de la
  méthode utilisée~(distance de Manhattan) posent un problème
  d'interprétation des déformations. %
  Les quatre cubes à droite symbolisent l'opération de séparations d'un
  espace-temps de cellules à quatre bits en quatre espace-temps de
  cellules à un bit. %
  La figure~\mygetdump{tsum-128x3-trans-bits} présente des
  transparences sur trois axes pour cinq \aquote{plans} de bits
  de~\arule{Qmean}, %
  La figure~\mygetdump{life-256x3-trans-bits} pour l'unique plan d'un
  espace-temps de~\arule{life}.}

\mydumpfull[htb]{anneal-256x3-diag}%
{}%
{Une séquence de tranches d'espace-temps de la règle~\arule{Anneal}}%
{Les premières coupes correspondent à un temps proche du temps initial,
  mais chaque tranche présente des caractéristiques archaïques plus ou
  moins marquées.}

\mydumpfull[htb]{tsum-128x3-faces}%
{}%
{Faces du cube d'espace-temps de~\arule{Qmean}}%
{Cette figure présente les six faces d'un cube d'espace-temps de la
  règle~\arule{Qmean}. %
  Les deux faces situées dans l'axe du temps, ainsi que leur âges
  sont clairement identifiables.}

\mydumpfull[p]{tsum-128x3-trans-bits}%
{}%
{Une transparence éclatée par plan de bits}%
{Les cinq bits de poids fort de la règle~\arule{Qmean} sont utilisés
  respectivement pour réaliser une
  transparence~\mygetfig[paren]{transcub} sur les trois axes de
  l'espace-temps. %
  Ces séquences sont à comparer à la
  figure~\mygetdump{tsum-128x3-trans-byte} utilisant la totalité des
  bits de chaque cellule pour la transparence.}

\mydumpfull[p]{tsum-128x3-trans-byte}%
{}%
{Une transparence utilisant tout l'espace de valeur d'une cellule}%
{Cette transparence fournit une image plus habituelle d'une
  dynamique cellulaire considérée comme un voxmap\footnote{%
    Le voxel est la version 3D du pixel.} %
  Des expériences antérieures menées avec des packages de visualisation
  scientifique \aquote{classiques}\footnote{%
    Au sens ou ils travaillent nécessairement en flottants, ou considèrent
    des données dont la localisation n'est pas implicite
    et en tout cas ne sont pas adaptés à un travail au niveau du bit. %
    Nous avons recherché, porté et expérimenté les outils suivants:~%
    \mysoftcite[XDataSlice]{XDS} et son format %
    \mysoftcite{HDF}, %
    \mysoftcite{VIS-5D}~\cite[]{Hibbard94}, %
    \mysoftcite{viewit} et %
    \mysoftcite{LIC}~\cite[]{Cabral93}.} %
  donnent pour ce genre de visualisation des résultats proches en terme
  d'images produites\footnote{%
    Mais pas en termes de ressources consommées pour la production. %
    Un autre inconvénient de ces packages, étant les erreurs
    introduites par la conversion flottante intermédiaire lorsque
    l'espace 2D de transparence produit est destiné à la réalisation
    de mesures statistiques, autant que d'images.} %
  de ceux présentés ici, %
  mais l'usage de nos fonc\-tions de manipulation d'espaces cellulaires
  est nécessaire\footnote{%
    D'un point de vue pratique, la conversion d'un bitmap
    tridimensionnel en voxmap de flottants, éventuellement reconverti
    en voxmap de bytes, est en effet parfaitement inadéquate.} %
  à la pro\-duc\-tion des transparences par plan de bits de
  la figure~\mygetdump{tsum-128x3-trans-bits}. %
  La comparaison des images produites par les méthodes plan de bits
  \vs voxel, suggère l'étude des relations entre les deux
  formulations, et une possible généralisation de la première à
  l'analyse de voxmaps.} %

\mydumpfull[htb]{anneal-256x3-trans-bits}%
{}%
{Une transparence par plan de bits pour~\arule{Anneal}}%
{La méthode de production commentée en
  figure~\mygetfig{transcub} est utilisée 
  pour la règle~\arule{Anneal}. %
  Le nombre d'états de~\arule{Anneal} étant égal à deux, la méthode
  éclatée se confond avec la méthode standard.}

\mydumpfull[htb]{life-256x3-trans-bits}%
{}%
{Une transparence par plan de bits pour~\arule{Life}}%
{La méthode de production commentée en
  figure~\mygetfig{transcub} est utilisée 
  pour la règle~\arule{Life}}

\mydumpfull[htb]{tsum-128x3-relmovedge}%
{}%
{Une détection de contour appliquée sur une projection relativiste}%
{L'accumulation de transformations renforce le problème de
  l'interprétation des images produites. %
  Ces figures illustrent les limites rencontrées en combinant une
  détection de contour une à projection relativiste à cellule de
  référence mobile, appliquée sur les trois axes de projection d'un
  cube d'espace-temps de la règle~{Qmean}.}

L'état actuel du développement du workbench n'a pas encore permis
l'expérimentation de constructions complexes d'objets mixtes. %
La mesure de variables d'états macroscopiques, par exemple, utilise
actuellement des outils externes, au lieu d'algorithmes SIMD adaptés à
la nature discrète et finie des espaces cellulaires. %

%%
%% Local Variables:
%% mode: auto-fill
%% iso-accents-minor-mode: nil
%% TeX-command-default: "mytex"
%% TeX-master: "top"
%% Local IspellParsing: "latex-mode ~latin1"
%% Local IspellPersDict: "~/.ispell_lisp"
%% LocalWords: "d\'evelop pement d'AC cellsim scamper CA LAB ca lab xlife hodge"
%% LocalWords: "HODGE-C cellang impl\'ementation d'espace-temps mips m\'egabytes"
%% LocalWords: "byte kilobyte m\'egaoctet gigaoctet bidimensionnel Moore kilobit"
%% LocalWords: "Eisenstein Gance Godard Marey Chronophotographe pixillation ur"
%% LocalWords: "McLaren Bergson Z\'enon d'\'El\'ee finitude ph\'enom\'enologique DX Linux"
%% LocalWords: "pip\'e dumpant dev null monochrome interfac\'ee Tcl Tool Command htb"
%% LocalWords: "Language s\'equencement casNew Workbench dumper Tk Anneal shell ad"
%% LocalWords: "s\'equenceur monochromes plotters hoc packages hodge-podge Gerhard"
%% LocalWords: "Schuster Beluso v-Z macro-g\'en\'eration habotinski workbench SIMD"
%% LocalWords: "SISD bottom-up lookup-tables manipulables artefactuelle NEW tsum"
%% LocalWords: "superluminique subluminique luminique timepyrcub relativiste rel"
%% LocalWords: "diag anneal Manhattan Bresenham Pendleton relmov relmovedge XDS"
%% LocalWords: "Qmean the-rules tsum-grey tsum-ray-one tsum-ray-all paren voxmap"
%% LocalWords: "est-ouest transcub trans-byte trans-bits voxel XDataSlice HDF"
%% LocalWords: "viewit LIC bitmap bytes versus voxmaps nord-sud name"
%% LocalWords: "lookup-table relativistes glider point-fixe"
%% LocalWords: "espace-temps l'espace-temps superglider d\'esaxage"
%% End:

%% Objectifs et contraintes
%% Outils existants
%% Calcul et visualisation

\section{Pourquoi notre nouvel environnement?}

Tout d'abord, m\^eme si l'histoire de l'\'etude des AC est presque aussi
ancienne que l'histoire des ordinateurs~\cite[]{Metropolis93b}, les
architectures cellulaires construites, bien que non encore
industrialis\'ees apparaissent comme les plus prometteuses\footnote{%
  Il est difficile de savoir si ce fait est \`a mettre au cr\'edit des
  architectures cellulaires, nous nous bornerons \`a proposer l'id\'ee que
  la perception du potentiel d'usage d'une architecture est
  inversement proportionnelle \`a sa r\'ealisation industrielle.}, %
et le probl\`eme des m\'ethodes de constructions d'architectures et
d'algorithmes cellulaires restent largement ouvert. %
De m\^eme, \`a propos du probl\`eme de
l'auto-reproduction~\cite[]{vonNeumann66}, la dualit\'e opposant
analytique \`a synth\'etique accompagne la naissance des AC\footnote{%
  Voir \cite{Reggia93} pour une approche permettant d'accorder plus
  d'autonomie aux ph\'enom\`enes d'auto-reproduction. %
  Nous consid\'erons, dans ce contexte, l'autonomie comme la mesure
  d'une position sur l'axe qui va du construit \`a l'\'emergeant.}. %
Cette dualit\'e, ou plut\^ot la capacit\'e offerte par les AC de basculer
d'un univers \'epist\'emologique \`a un autre tout en restant au sein d'un
mod\`ele unique, doit \^etre incarn\'ee dans l'environnement. %
Nous souhaitons donc un environnement permettant une approche mixte. %

Par ailleurs, aucun des outils que nous avons cit\'es ne r\'epond
totalement \`a un ensemble de contraintes plus pratiques que nous avons
fix\'ees comme cadre de d\'eveloppement. %
\begin{itemize}
\item la nature publique\footnote{%
    Par \aconcept{publique} nous entendons accessible sous forme
    source et utilisable sans licence commerciale.} %
  des outils
\item la rapidit\'e d'ex\'ecution
\item la richesse de visualisation
\item la souplesse de sp\'ecification
\item un environnement d'exploitation Unix et X11
\item la modularit\'e des constituants en vue de la construction et du
  s\'equencement d'exp\'eriences cellulaires. %
\end{itemize}

\asoft{CA~LAB} n'est pas un outil public et n'est pas disponible dans
l'environnement Unix. %
\asoft{Cellsim} de par sa d\'ependance \`a~\asoft{SunView} n'est
utilisable que dans un environnement Sun\footnote{%
  Felicity George a r\'ealis\'e en 1992 un portage de~\asoft{cellsim} pour
  l'environnement~\asoft{XView}.}, %
\asoft{scamper} pr\'esente une application int\'egrant un formalisme de
sp\'ecification des r\`egles bas\'ees par pattern, des compteurs de
populations et une interface visuelle bas\'ee sur X11, mais il ne permet ni
ex\'ecution rapide, ni s\'equencement. %
\asoft{Cellang} propose un langage de description d'AC, mais son
compilateur ne produit pas un code permettant de tirer le meilleur
parti des ressources disponibles. %

\section{Des AC pour l'optimisation d'usage de ressources}

Concernant le parall\'elisme massif, les tendances de la pr\'ec\'edente
d\'ecennie en architecture des machines ne confirment pas les
pr\'edictions avanc\'ees \`a son propos\footnote{%
  Machines SIMD massives reposant sur des techniques de fabrication de
  type wafer scale int\'egration, associ\'ees \`a des algorithmes de
  reconfigurations de r\'eseaux d'interconnexion permettant la
  r\'esistance aux d\'efauts; pyramides de telles architectures pour la
  vision et plus g\'en\'eralement l'IA, ou architectures diverses pour la
  physique, la biologie mol\'eculaire, le pattern matching haut d\'ebit,
  la r\'ealisation de film holographique temps r\'eel.}. %
N\'eanmoins, hors consid\'erations architecturales, les ressources
accessibles continuent \`a cro\^itre \`a un rythme soutenu\footnote{%
  Au point qu'on pourrait mettre en doute l'application du
  principe des gaz parfaits \`a la consommation de ressources. %
  Contre cette tendance, nous proposons de consid\'erer qu'il n'est plus
  temps d'apprendre comment se contenter de ressources restreintes,
  mais comment tirer profit de ressources disponibles.}. %
Hormis le calcul intensif classique, les gains de ressources semblent
s'orienter vers le parall\'elisme MIMD massif pour la gestion et le
support multim\'edia pour le poste de travail. %
Nous affirmons que les AC constituent un mod\`ele efficace et universel
pour l'utilisation de ces ressources, y compris par la prise en compte
du parall\'elisme multi-calculateur\footnote{%
  C'est \`a dire l'utilisation d'un ensemble de machines
  autonomes~(stations de travail ou autres, localement parall\`eles ou
  non) reli\'ees par un r\'eseau d'interconnexion~(de type bus ou point
  \`a point) dot\'ees d'une couche logicielle syst\`eme ad hoc, comme une
  ressource parall\`ele en soi.} %
associ\'e \`a l'interconnexion par \astrangeterm{WAN}{%
  Wide Area Network, il s'agit ici de remarquer que les d\'ebits
  disponibles ou annonc\'es pour les r\'eseaux de t\'el\'ecommunication publics
  permettent d'envisager la construction virtuelle de
  machines cellulaires reparties sur une large zone g\'eographique.} %
de ceux-ci. %
La r\'egularit\'e structurelle des donn\'ees et des traitements des AC
associ\'es \`a l'absence de probl\`eme algorithmique de synchronisation
des \'el\'ements du calcul r\'eparti, permet un calcul et une affectation de
ressource simple. %
Une approche bottom-up de la sp\'ecification du comportement des
\'el\'ements cellulaires permet, de plus, de s'affranchir des probl\`emes
d'architecture CPU, une fonction consid\'er\'ee sous l'angle de sa table
de transition est ind\'ependante des architectures de
traitement\footnote{%
  A une permutation de l'ordre des bits sur le mot pr\`es, et toujours en
  consid\'erant que l'unit\'e de traitement r\'eellement disponible reste
  fondamentalement une architecture de von Neumann dot\'ee d'une m\'emoire
  adressable importante.}. %

\mysetref{repartition-du-calcul}
A propos de la r\'epartition d'un calcul cellulaire sur r\'eseau de
stations de travail\footnote{%
  Algorithme distribu\'e avec recouvrement: calcul des bords, envoi des
  bords, calcul du centre, r\'eception des bords.}, %
et concernant le probl\`eme de l'overhead associ\'e au passage de messages
impl\'ementant la synchronisation des bords lorsque le calcul courant
r\'eclame un acc\`es \`a une cellule voisine, ind\'ependamment d'une mesure
particuli\`ere des d\'ebits respectifs du couple
\bialt{cpu}{m\'emoire}~($D_{m}$) et du canal de communication($D_{c}$),
le rapport des tailles \bialt{locales}{distantes} $N^{d}/N^{d-1} = N$
garantit l'existence d'une valeur $k$ de $N$ telle que l'usage de
plusieurs n{\oe}uds de calculs sera plus efficace que l'usage d'un
seul d'entre eux:~$\frac{D_{m}}{D_{c}} k > 1$. %
%% \[\frac{D_{m}k^{d}}{D_{c}k^{d-1}} > 1\]

A d\'efaut de machines cellulaires industrielles, nous
esp\'erons pouvoir utiliser les DSP et les acc\'el\'erateurs
\aquote{bitblt}~\cite[]{Pike84} aujourd'hui largement disponibles
comme acc\'el\'erateurs de nos moteurs cellulaires. %

\section[AC et morphog\'en\`ese]{Des AC comme outils de r\'eflexion morphog\'en\'etique}

L'essence de la forme pose un probl\`eme, ne serait ce que
par l'existence de la g\'en\'ericit\'e du moule. %
Tout \'etat du monde n'est pas une forme, toutefois une forme est par
d\'efinition une configuration d'espace cellulaire prise parmi un
ensemble fini de configurations. %
Cette d\'efinition r\'eclame un crit\`ere de
\bialt{discrimination}{r\'eduction} de l'ensemble des \'etats. %
Une premi\`ere possibilit\'e r\'eside dans l'utilisation d'une fonction
arbitraire sur un \'etat d\'eterminant la valeur d'un attribut formel:
esth\'etique, g\'eom\'etrique, sta\-tistique, arithm\'etique. %
Les formes sont figuratives, abstraites, s\'eduisantes, rondes, carr\'ees,
sym\'etriques, rares, uniques\footnote{%
  Chaque configuration d'espace est \'evidemment unique, on pourrait
  n\'eanmoins mesurer, intuitivement ou formellement, l'unicit\'e \`a partir
  du produit de l'ensemble des \'etats par l'ensemble des
  fonctions de qualification formelle disponibles. %
  Les fonctions de qualification sont des fonctions de r\'eduction de
  l'information~(le nombre de bits rendu est inf\'erieur au nombre de
  bits fourni), le cardinal de l'image d'un ensemble de formes est
  donc inf\'erieur \`a celui de l'ensemble de formes, la distribution des
  redondances de valeurs induites par la diff\'erence de cardinalit\'e
  d\'eterminant les classes d'\'equivalence de cet attribut. %
  
  Une relation d'ordre pour l'unicit\'e est alors d\'efinie par une
  fonction minimisant le cardinal de l'union des classes d'\'equivalence
  des images des formes.}, %
r\'eguli\`eres, fractales, denses, color\'ees, leur num\'ero d'ordre enfin,
peut \^etre argument d'une fonction choisie au hasard. %
Cette m\'ethode pose le probl\`eme de l'\'enum\'eration des fonctions de
r\'eduction\footnote{%
  Ainsi que celui du choix \`a priori des \'etats consid\'er\'es. %
  Les fonctions de qualification sont implicitement destin\'ees \`a
  introduire une r\'eduction du nombre des formes par l'introduction de
  classes d'\'equivalence sur l'ensemble des \'etats.}. %
Une deuxi\`eme m\'ethode consiste \`a opposer la g\'en\'eticit\'e \`a la
g\'en\'ericit\'e, c'est \`a dire l'histoire de la forme au moule de la forme. %
Dans ce contexte la r\'eduction du nombre d'espaces appel\'es \`a la forme
est d\'etermin\'ee par l'existence d'une trajectoire dans l'espace des
\'etats. %
Cette trajectoire \'etant elle m\^eme d\'etermin\'ee par l'application it\'er\'ee
d'une fonction de l'ensemble des \'etats vers l'ensemble des \'etats. %
Les trajectoires ne sont pas moins nombreuses que les \'etats, mais si
l'on ajoute une contrainte de localit\'e aux fonctions g\'en\'eratrices de
trajectoires, une tentative d'\'enum\'eration devient accessible. %
Cette contrainte de localit\'e constitue le fondement d'une approche
cellulaire de la forme. %
Cette contrainte introduit un lien entre la vitesse et la forme. %
Plus g\'en\'eralement, la notion de contrainte s'applique ais\'ement au
monde des formes dans une perspective esth\'etique, le libre choix des
contraintes de production de l'objet d\'eterminant la qualit\'e artistique
de l'{\oe}uvre, dont la forme peut \^etre con\c{c}ue comme une trace
d'ex\'ecution\footnote{%
  La perception de l'{\oe}uvre n'est jamais exempte de la perception
  de l'histoire de sa fabrication~\cite[]{Gombrich71}.}. %
De m\^eme la \aquote{forme} d'un espace cellulaire
devient la trace d'un programme\footnote{%
  De l'application d'une loi non ergodique sur un espace initial
  n\'ecessairement al\'eatoire.}. %

\section[AC, espace, temps et calcul]%
{Des AC comme outils de r\'eflexion sur l'espace, le temps et le calcul}

Nous avons d\'ej\`a mentionn\'e l'existence d'une dualit\'e
analytique-synth\'etique. %
D'autres projections de cette dualit\'e sont possibles, sur les axes
dynamique statique, heuristique algorithmique, parall\`ele s\'equentiel ou
temporel spatial. %
La trajectoire d'un AC sur un espace de configuration est un objet
dynamique, la structure d'un espace de r\`egles de transition d'AC pour
un voisinage donn\'e est un objet statique. %
A la fronti\`ere se trouvent des objets donn\'es par l'application
it\'erative d'une fonction simple sur un espace cellulaire initial
globalement param\'etr\'e~(l'automate auto-reproducteur de von Neuman, ou
une machine de Turing impl\'ement\'ee par un espace de configuration
d'un AC tel \arule{Life} de Conway). %

La quantit\'e de parall\'elisme d'une ex\'ecution peut se mesurer au nombre
d'instructions diff\'erentes rencontr\'ees lors de cette ex\'ecution. %
Il existe une \'equivalence entre un AC d\'efini par un r\'eseau
d'interconnexion et une fonction de transition, et un algorithme
cellulaire d\'efinit par une s\'equence d'instructions
cellulaires. %
Si un AC est consid\'er\'e comme une instruction cellulaire \'el\'ementaire,
celle-ci poss\`ede un co\^ut, qui s'exprime comme la taille de la table
de v\'erit\'e de cette instruction cellulaire. %
Si un AC est consid\'er\'e comme une s\'equence d'instructions cellulaires,
celle-ci poss\`ede un co\^ut, qui s'exprime comme la somme des temps de
r\'ealisation de chaque instruction, et in fine la somme de la taille
des tables de v\'erit\'e de chaque instruction \'el\'ementaire. %
L'impl\'ementation optimale d'un algorithme cellulaire serait obtenu
lorsque, pour une occupation maximale de la m\'emoire disponible pour les
tables de v\'erit\'es~(les fonctions de transition) des instructions
cellulaires, le nombre de primitives diff\'erentes composant l'algorithme
est minimum. %
Le ratio du nombre de primitives employ\'ees sur le nombre de primitives
existantes d\'etermine l'importance des effets relevant de la
dynamique~(l'heuristique) de ceux relevant de l'algorithme proprement
dit dans la fonction associ\'ee au programme cellulaire.

Nous disons que la fonction de transition d'un AC est localement
univer\-selle\footnote{%
  Nous dirons parfois qu'un AC est globalement universel, pour
  d\'esigner la propri\'et\'e usuelle d'universalit\'e par opposition \`a la
  propri\'et\'e d'universalit\'e locale des PE de machines SIMD\@. %
  L'universalit\'e locale implique \'evidement l'universalit\'e globale. %
  \arule{Life} constitue dans le contexte de ce travail l'arch\'etype
  d'un AC dot\'e de l'universalit\'e globale, et le processeur GAPP
  l'arch\'etype d'un AC dot\'e de l'universalit\'e locale.} %
si une seule cellule consid\'er\'ee comme PE se comporte comme une machine
universelle sous r\'eserve d'une extension virtuelle du nombre d'\'etats
cellulaires. %
C'est a priori le cas de toutes les architectures SIMD,
m\^eme si le faible nombre d'\'etats accessibles\footnote{%
  Ce faible nombre d'\'etats doit n\'eanmoins \^etre exprim\'e en logarithme,
  c'est \`a dire en nombre de bits adresse.} %
destine \'evidemment ces machines \`a un usage cellulaire. %

D'un cot\'e le \aquote{calcul s\'equentiel} pur d\'eroule une s\'equence
d'instructions d\'efinie en extension sur un espace initial \'egalement
d\'efini en extension, de l'autre le \aquote{calcul cellulaire pur}
it\`ere une loi unique d\'efinie en extension sur un espace initial dont
seules des propri\'et\'es statistiques\footnote{%
  La taille d'un espace cellulaire incarn\'e, n\'ecessairement fini, doit
  n\'eanmoins \^etre consid\'er\'ee comme tendant vers l'infini; la
  sp\'ecification d'une configuration cellulaire particuli\`ere en
  extension devient alors une barri\`ere \`a la pleine utilisation des
  ressources disponibles.} %
sont d\'efinies. %
La distinction entre une s\'equence d'instructions et une \emph{loi}
apparaissant comme une instruction unique, ne tient ni \`a la taille de
la loi ni \`a une propri\'et\'e telle l'universalit\'e de l'architecture
sous-tendant la s\'equence d'instructions, ou l'universalit\'e
locale ou globale de la loi\footnote{%
  En effet toute s\'equence d'instruction, tout algorithme peuvent
  \'evidemment \^etre transform\'es en une instruction
  unique par le calcul d'une lookup-table.} %
mais \`a l'opposition entre l'unicit\'e de l'application d'une s\'equence
d'instructions, d'un algorithme et la n\'ecessaire it\'eration sans
fin\footnote{%
  Comme pour l'espace, ressource n\'ecessairement limit\'ee devant
  n\'eanmoins \^etre consid\'er\'ee comme tendant vers l'infini, le temps
  cellulaire, le nombre de cycles d'it\'erations de l'application de la
  loi sur l'espace doit \^etre consid\'er\'e comme tendant vers l'infini.} %
de la loi sur l'espace cellulaire. %
A la fronti\`ere\footnote{%
  La structure de l'espace englobant le s\'equentiel et le cellulaire
  appara\^it ici tr\`es clairement.} %
le calcul cellulaire it\'erant une loi unique sur un espace dont la
configuration ini\-tiale est une transformation~(un mapping) d'une s\'erie
d'instructions d'un calcul s\'equentiel, c'est \`a dire le calcul
cellulaire exploitant l'universalit\'e globale d'une loi; ou encore le
calcul cellulaire SIMD d\'eroulant une s\'equence d'instructions d\'efinie
en extension, c'est \`a dire le calcul cellulaire utilisant
l'universalit\'e locale d'une loi, et sans doute, mais sans qu'il soit
possible d'en \'etablir simplement la part, l'universalit\'e globale de
cette loi. %

Si l'on consid\`ere que l'encodage d'un algorithme sur une machine
s\'equentielle pose un probl\`eme aigu, et encore davantage
l'encodage du probl\`eme \`a traiter, il en va bien s\^ur de m\^eme pour les
machines SIMD. %
On continue toutefois dans le contexte SIMD \`a r\'esoudre des probl\`emes
algorithmiques, \`a utiliser l'universalit\'e pour la construction de
machines algorithmes. %
Lorsque l'on passe \`a une machine cellulaire pure, la nature
intrins\`equement statistique du processus de calcul\footnote{%
  Nous avons jusque l\`a \'evoqu\'e la seule incertitude statistique
  associ\'ee \`a la sp\'ecification de la configuration initiale, la loi
  \'etant consid\'er\'ee comme purement d\'eterministe, on voit bien cependant
  comment un ind\'eterminisme sur l'espace est de fait indiff\'erentiable
  d'un ind\'eterminisme sur la loi, laquelle n'est jamais que de
  l'espace local g\'en\'erateur du temps.} %
change radicalement le point de vue. %
Le probl\`eme \`a traiter ne pouvant d\`es lors gu\`ere s'exprimer que par
analogie \`a une morphog\'en\`ese
cellulaire\footnote{%
  Expression \`a prendre dans un sens non restrictif \`a la biologie bien
  s\^ur, \`a l'existence pr\`es d'un paysage \'epig\'en\'etique.}. %
Il sera toujours possible de plonger une approche SIMD, voire
purement s\'equentielle et distribu\'ee dans un contexte cellulaire pur,
mais outre que cela pose des probl\`emes qui nous apparaissent assez
ardus de faisabilit\'e pratique\footnote{%
  Cela pose aussi des probl\`emes plus fondamentaux, dont le moindre
  n'est pas d'\'eclaircir la nature ontologique de l'incertitude
  statistique qui semble bien devoir \^etre li\'ee \`a une d\'efinition d'un
  calcul cellulaire r\'eellement massif, m\^eme si la dimension
  technologique de cette incertitude existe \'egalement.}, %
cela ne tend pas \`a r\'epondre \`a la
question de la sp\'ecificit\'e du calcul cellulaire pur. %
En effet, si celui-ci ne semble pas, dans son approche la plus stricte,
isomorphe \`a une machine universelle arbitraire, ce pourrait \^etre d\^u \`a
la nature particuli\`ere des morphog\'en\`eses cellulaires, lesquelles ne
sont peut-\^etre pas r\'eductibles \`a des algorithmes. %
Le calcul s\'equentiel peut \^etre consid\'er\'e comme une version
adimensionnelle~(de dimension z\'ero) d'un calcul cellulaire. %
Le calcul cellulaire algorithmique\footnote{%
  C'est \`a dire la r\'ealisation d'un algorithme par une une
  configuration cellulaire initiale associ\'ee \`a une loi globalement
  universelle.} %
correspondant au choix d'un point unique dans l'espace des \'etats
cellulaires, permettant par le plongement d'une machine de Turing dans
l'espace cellulaire, de r\'ealiser une r\'eduction
dimensionnelle de l'espace cellulaire. %
Le calcul cellulaire SIMD\footnote{%
  C'est \`a dire la r\'ealisation d'un algorithme sur une configuration
  cellulaire par une m\'ethode s\'equentielle locale.} %
correspondant au choix d'un mapping probl\`eme vers configuration cellulaire,
associ\'e ou non \`a l'utilisation des propri\'et\'es dynamiques d'une loi ou
plusieurs lois pour contr\^oler une trajectoire dans l'espace des
configurations cellulaires. %
Le calcul cellulaire pur, par la prise en compte de l'incertitude
statistique, cherche \`a dresser un paysage g\'en\'eral du calcul,
lequel devient une propri\'et\'e intrins\`eque, voire \'emergeante d'un mod\`ele
d'univers compos\'e d'atomes d'espace dot\'e d'une loi universelle. %
Une des propri\'et\'es particuli\`erement int\'eressante, de notre point de
vue, du calcul cellulaire, r\'eside dans les \'equivalences espace-temps
rendues possibles par son architecture. %
En effet la mesure des quantit\'es \'el\'ementaires de cette architecture,
ainsi que les transformations mutuelles possibles entre ces quantit\'es
sont d'impl\'ementation ais\'ee. %
Ainsi, en calcul cellulaire pur, le nombre d'\'etats d'une cellule et le
nombre de voisins de cette cellule d\'etermine une borne sup\'erieure \`a la
complexit\'e\footnote{%
  Nous n'avons pas de d\'efinition pr\'ecise de cette complexit\'e,
  il s'agit a priori d'une mesure statistique de la richesse de
  comportement, rapport\'ee \`a une dur\'ee. %
  La capacit\'e et la rapidit\'e morphog\'en\'etique associ\'ees \`a une r\`egle.} %
du calcul effectu\'e\footnote{%
  L'analyse de la table de transition permet de pr\'eciser une valeur
  de complexit\'e, en particulier la distance \`a la borne sup\'erieure. %
  Une table constante, les tables identit\'e ou inverse sur un seul
  bit du mot de voisinage~(translation, et translation alternance de
  la configuration initiale) sont faciles \`a identifier. %
  Une g\'en\'eralisation possible de ces notions conduit \`a la d\'efinition
  de param\`etres de contr\^ole de l'espace des r\`egles tels que le
  param\`etre $\lambda$ de Langton ou le param\`etre d'ench\^assement.}. %
La taille de \arule{Life} est de~512~bits,
celle de l'automate auto-reproducteur de von Neumann de
$2^{5\lceil\log_{2}(29)\rceil}\lceil\log_{2}(29)\rceil =
20 \text{m\'egaoctets}$. %
Il serait int\'eressant de conna\^itre la taille minimale d'une
configuration initiale de \arule{Life} r\'ealisant une machine
universelle, il n'existe pas \`a notre connaissance de r\'ealisation
effective d'une telle configuration. %
La taille de la configuration r\'ealisant l'automate auto-reproducteur
de von Neumann a \'et\'e \'evalu\'ee \`a plusieurs centaines de milliers de
cellules. %
%% John kemeny 1955
Jaqueline Signorini a \'etudi\'e les conditions de r\'ealisation de cet
automate sur le GAPP~\cite[]{Signorini92}, mais une configuration
auto-reproductrice compl\`ete n'a encore jamais \'et\'e men\'ee \`a bien. %

L'impl\'ementation de \arule{Life} sur une machine s\'equentielle
$M_{\text{SISD}}$ utilise $a$ cycles de base pour chacune des $N$
cellules, et $b = 2$ bits\footnote{% 
  Un plan \avar{past} et un plan \avar{futur}, impl\'ementant le
  synchronisme de mise \`a jour, on peut envisager des optimisations
  r\'eduisant le nombre de bits n\'ecessaire au plan secondaire,
  lesquelles se ram\`enent \`a une compression de la configuration
  cellulaire, associ\'ee \`a des m\'ethodes d'acc\`es sp\'ecialis\'ees au codage de
  la compression, renfor\c{c}ant ainsi la tendance des machines
  s\'equentielles \`a l'\'echange de temps contre de l'espace.} %
par cellule pour le calcul d'une g\'en\'eration de \arule{Life}, plus
$c \times  i$ bits pour le codage des instructions, $i$ \'etant la taille
moyenne d'une instruction. %
$c$ et $a$ sont typiquement tous les deux de l'ordre de quelques
centaines d'\bialt{instructions}{cycle} \'el\'ementaire d'une machine
monoprocesseur id\'eale pour une version na\"ive~(non optimis\'ee) de
l'\'emulateur de la machine cellulaire. %
Une impl\'ementation par lookup-table avec collecte pipelin\'ee du
voisinage augmente $c$ et diminue $a$ d'un ordre de grandeur. %

\begin{equation}
  \frac{\text{Space}}{\text{Time}} = \frac{bN+ci+M_{\text{SISD}}}{aN}
  \label{eq-life-sisd}
\end{equation}

Avec:\[b = 2,~i = 32,~c \approx a \approx 10^2,~M_{\text{SISD}}
\approx 10^6\]

La valeur de $a$ et $c$ d\'ependent du nombre $k$ d'\'etats par cellule et
du rayon $r$ du voisinage cellulaire. %
La taille du mot de voisinage~(9~dans le cas de \arule{Life})
d\'etermine ces valeurs. %

L'impl\'ementation de~\arule{Life} sur une machine
SIMD~$M_{\text{SIMD}}$ utilise $a^{\prime}$~cycles de base pour le
calcul d'une g\'en\'eration, ind\'ependamment du nombre de cellules,
et~$b^{\prime}$~bits par cellule, les~$c^{\prime} \times  i^{\prime}$~bits
d'instructions \'etant traditionnellement consid\'er\'es comme n'appartenant
pas stricto sensu \`a la machine SIMD, mais \`a une indispensable machine
s\'equentielle, h\^ote de toute machine SIMD\@. %
Les $a^{\prime}$~cycles de calcul d\'ependent de la machine SIMD choisie,
une machine SIMD typique est plus difficile \`a proposer qu'une
machine SISD id\'eale, sur le GAPP, de l'ordre de la centaine pour une
impl\'ementation a priori de la fonction de transition\footnote{%
  L'algorithme arithm\'etique de \arule{Life}, comptage plus seuil.}, %
de l'ordre du millier pour une impl\'ementation a posteriori de celle-ci\footnote{%
  Une d\'erivation algorithmique du DFA de \arule{Life} calcul\'ee \`a partir
  de la table de transition, laquelle peut \^etre obtenue par observation
  d'une trajectoire sur l'espace des configurations, associ\'ee \`a une
  algorithmique SIMD g\'en\'erique interpr\'etant le DFA sur le mot de voisinage.}. %
De l'ordre d'une dizaine de bits pour les $b^{\prime}$ bits d'espace. %

\begin{equation}
  \frac{\text{Space}}{\text{Time}} = \frac{(b^{\prime}+M_{\text{SIMD}})N}{a^{\prime}}
  \label{eq-life-simd}
\end{equation}

Avec:\[b \approx 10,a^{\prime} \approx 10^2,~\text{SIMD} \approx 10^2\]

L'impl\'ementation de \arule{Life} sur une machine cellulaire pure $M_{\text{CA}}$
utilise $a^{\prime\prime} = 1$ cycle et  $b^{\prime\prime} = 1$ bit par
cellule. %
Par contre le co\^ut li\'e \`a l'architecture s'exprime simplement comme la
taille de la table de transition. %

\begin{equation}
  \frac{\text{Space}}{\text{Time}} =
  \frac{(b^{\prime\prime}+M_{\text{CA}})N}{a^{\prime\prime}}
  \label{eq-life-cell}
\end{equation}

Avec:\[a^{\prime\prime} = 1,~b^{\prime\prime} = 1,~\text{CA} = 512\]

Les valeurs d'espace propos\'ees pour le terme $M_{\text{SISD}}$ et le
facteur $M_{\text{SIMD}}$ dans les \'equations~\ref{eq-life-sisd}
et~\ref{eq-life-simd} repr\'esentent un nombre de transistors CMOS id\'eal
pour ces architectures. %
Le facteur $M_{\text{CA}}$ de l'\'equation~\ref{eq-life-cell} n'est pas
un nombre de transistors approxim\'e, mais un nombre exact de bits. %
Cette diff\'erence de type entre les objets compt\'es se trouve
consid\'erablement pond\'er\'ee par les croissances respectives des tailles
d'espaces n\'ecessaires \`a la r\'ealisation d'une fonction de transition
lorsque la fonction de transition devient un param\`etre des \'equations
pr\'ec\'edentes. %
Sur une machine SISD l'augmentation de la taille du mot de voisinage
induit une croissance lin\'eaire de la taille de l'espace et du temps
utilis\'e pour calculer une trajectoire cellulaire. %
L'augmentation de la taille d'espace ne d\'epend que de la taille du mot
de voisinage, pas de la complexit\'e de la fonction de
transition\footnote{%
  A condition, bien s\^ur, de ne pas transformer la fonction de
  transition en lookup-table.}, %
mais l'augmentation du temps de calcul d\'epend de la taille de la
configuration initiale. %
Sur une machine SIMD l'augmentation de la taille du mot de voisinage
induit \'egalement une croissance lin\'eaire de la taille de l'espace et
du temps utilis\'e pour calculer une trajectoire cellulaire. %
Par contre l'augmentation de la taille d'espace d\'epend de la taille
du mot de voisinage et de la complexit\'e de la fonction de transition,
ce facteur de croissance de la taille d'espace est donc sup\'erieur \`a
celui d'une machine SISD\@. %
De plus la taille de l'architecture est fonction de la taille de la
configuration initiale, ce facteur pr\'esentant \'evidemment la source
principale des besoins de ressource spatiale. %
Par contre, pour prix de cette augmentation lin\'eaire des besoins
d'espace par rapport \`a la version SISD, le temps de calcul n'est plus
fonction de la taille de l'espace initial. %
Sur une machine cellulaire pure l'augmentation de la taille du mot de
voisinage induit une croissance exponentielle de la taille d'espace
via une croissance exponentielle de la taille de la table elle-m\^eme
facteur de la taille de la configuration initiale, le temps se
r\'eduisant \`a l'unit\'e. %

\section{Plan de lecture}

Dans le chapitre~\mygetref{space-time-law} nous pr\'esentons la
structure des objets de notre \'etude:~l'espace, le temps et les lois, %
ainsi que les m\'ethodes applicables sur ces objets:~\'edition, r\'eduction,
transformation, visualisation pour l'espace et le temps, plus
construction, alt\'eration, composition, m\'elange et s\'eparation pour les
lois. %
Nous pr\'esentons \'egalement dans ce chapitre l'architecture de nos
moteurs cellulaires. %

Dans le chapitre~\mygetref{rule-space} nous pr\'esentons les outils
conceptuels utilis\'es pour la quantification de l'espace des r\`egles, et
nous proposons un nouveau param\`etre de contr\^ole pour l'exploration et
la structuration de cet espace:~le param\`etre d'ench\^assement. %

Le chapitre~\mygetref{about-anneal} est consacr\'e \`a une \'etude
exp\'erimentale de la r\`egle~\arule{Anneal} bas\'ee sur l'exploitation de
nos outils. %
Les exp\'eriences r\'ealis\'ees nous am\`enent \`a proposer l'existence d'une
relative invariance spatio-temporelle des trajec\-toires associ\'ees \`a
une transition de phase dans l'espace des configurations initiales. %

La conclusion introduit les limites de notre approche qui \`a associ\'e
un constructivisme bottom-up \`a une \aquote{gestion de tables
  g\'en\'eralis\'ee}. %
Nous \'evoquons pour finir une approche alternative de r\'ealisation
cellulaire purement algorithmique. %

\mydumpfull{anneal-sum-1k-8k}%
{Une superposition de r\'eduction int\'egrale}%
{Une superposition de r\'eduction int\'egrale de \arule{Anneal}}%
{Cette image est obtenue par la superposition de la r\'eduction int\'egrale
  $\sum_{i=0}^{i=n}t_{i}xy$ des~$2^{10}$ et~$2^{13}$ premi\`eres
  g\'en\'erations d'un espace initial de taille~\asquare{1024} et de
  densit\'e\,\Sundemi.}

\mydumpfull{anneal-sum-1k-8k-xv}%
{}%
{Une autre colormap pour la figure
  \protect\mygetdump{anneal-sum-1k-8k}}%
{Une colormap en bandes. %
  Pour cette image, comme pour toutes les images en niveaux de gris
  que nous avons pr\'esent\'ees, un p\'eriph\'erique de sortie capable de
  restitution nuanc\'ee nous fait d\'efaut.}

%%
%% Local Variables:
%% mode: auto-fill
%% iso-accents-minor-mode: nil
%% TeX-command-default: "mytex"
%% TeX-master: "top"
%% Local IspellParsing: "latex-mode ~latin1"
%% Local IspellPersDict: "~/.ispell_lisp"
%% LocalWords: "l'auto-reproduction s\'equencement CA LAB cellsim SunView XView"
%% LocalWords: "Felicity scamper pattern cellang d'AC SIMD wafer scale l'IA MIMD"
%% LocalWords: "matching holographique multim\'edia multi-calculateur ad hoc WAN"
%% LocalWords: "Wide Area Network t\'el\'ecommunication bottom-up von Neumann cpu PE"
%% LocalWords: "adressable repartition-du-calcul l'overhead uds NEW DSP bitblt"
%% LocalWords: "morphog\'en\'etique g\'en\'ericit\'e cardinalit\'e fractales g\'en\'eticit\'e uvre"
%% LocalWords: "it\'er\'ee ergodique analytique-synth\'etique param\'etr\'e Neuman Turing"
%% LocalWords: "auto-reproducteur impl\'ement\'ee Conway univer GAPP it\`ere it\'erant"
%% LocalWords: "lookup-table mapping ind\'eterminisme indiff\'erentiable Langton DFA"
%% LocalWords: "\'epig\'en\'etique morphog\'en\`eses espace-temps Jaqueline Signorini past"
%% LocalWords: "auto-reproductrice monoprocesseur l'\'emulateur impl\'ementation"
%% LocalWords: "stricto sensu SISD CMOS approxim\'e space-time-law rule-space in"
%% LocalWords: "about-anneal Anneal constructivisme"
%% End:

%% Pourquoi notre nouvel environnement?
%% Optimisation d'usage de ressources
%% r\'eflexion morphog\'en\'etique
%% espace, temps et calcul
%% Plan de lecture

%%
%% Local Variables:
%% mode: auto-fill
%% iso-accents-minor-mode: nil
%% TeX-command-default: "mytex"
%% TeX-master: "top"
%% Local IspellParsing: "latex-mode ~latin1"
%% Local IspellPersDict: "~/.ispell_lisp"
%% LocalWords: "impl\'ementatoire SIMD Lattice Gas Automata cytosquelette"
%% End:
